	\section*{Abstract}
	The T0-Model offers two mathematically equivalent but conceptually different calculation methods for particle masses: the direct geometric method and the extended Yukawa method. Both approaches work with the single geometric constant $\xipar = \frac{4}{3} \times 10^{-4}$; the rational quantum numbers $(r_i, p_i)$ are structural choices, and $v$ serves as a conversion anchor. This documentation also contains neutrino quantum numbers and the quantum field theoretical derivation of the $\xi$-constant through EFT matching and 1-loop calculations. The neutrino part with double $\xi$-suppression is not tenable according to R109 (see the addendum at the end); for neutrinos, Doc.~340 is authoritative. For charged leptons and quarks, the average deviation is below 1\%; the over twenty free Standard Model mass parameters are thereby traced back to the single parameter $\xi$ and fixed rational quantum numbers.

	
	\section{Introduction}
	\label{sec:introduction}
	
	Particle physics faces a fundamental problem: The Standard Model with its over twenty free parameters offers no explanation for the observed particle masses. They are measured as inputs, not derived from a deeper principle. The T0-Model addresses this with two complementary calculation methods that work with the single parameter $\xi$ and rational quantum numbers; the originally included treatment of neutrino masses is not tenable according to R109 (see addendum).
	
	\subsection{The Parameter Problem of the Standard Model}
	\label{subsec:parameter_problem}
	
	Despite its experimental success, the Standard Model has a theoretical gap: It contains more than 20 free parameters that must be determined experimentally. These include:
	
	\begin{itemize}
		\item \textbf{Fermion Masses}: 9 charged lepton and quark masses
		\item \textbf{Neutrino Masses}: 3 neutrino mass eigenvalues
		\item \textbf{Mixing Parameters}: 4 CKM and 4 PMNS matrix elements
		\item \textbf{Gauge Couplings}: 3 fundamental coupling constants
		\item \textbf{Higgs Parameters}: Vacuum expectation value and self-coupling
		\item \textbf{QCD Parameters}: Strong CP phase and others
	\end{itemize}
	
	\begin{important}[Parameter Reduction]
		The T0-Model traces the over twenty free parameters of the Standard Model back to \textbf{one}: the geometric constant $\xipar = \frac{4}{3} \times 10^{-4}$, motivated by the geometry of three-dimensional space. In addition there are the rational quantum numbers $(r_i, p_i)$ as structural choices and $v$ as a conversion anchor. This version also contains neutrino quantum numbers (not tenable according to R109, see addendum) as well as the quantum field theoretical derivation.
	\end{important}
	
	\section{Methodological Clarification: Establishment vs. Prediction}
	\label{sec:methodological_clarification}
	
	\begin{important}[Scientific-Historical Classification]
		The T0-Model draws on the methodology of \textbf{pattern recognition and systematic classification}, which was successful, for example, for the periodic table (Mendeleev 1869) or the quark model (Gell-Mann 1964).
	\end{important}
	
	\subsection{Two-Phase Development}
	\label{subsec:two_phases}
	
	\textbf{Phase 1: Establishment of Systematics}
	\begin{enumerate}
		\item Pattern recognition in known particle masses (electron, muon, tau)
		\item Parameter determination from experimental data
		\item Establish quantum number assignment
		\item Demonstrate mathematical equivalence of both methods
	\end{enumerate}
	
	\textbf{Phase 2: Unfolding Predictive Power}
	\begin{enumerate}
		\item Extrapolation to unknown particles
		\item Derive quark sector from lepton patterns
		\item Predict new generations
		\item Conduct experimental tests
	\end{enumerate}
	
	\subsection{Historical Precedence of Successful Pattern Physics}
	\label{subsec:historical_precedence}
	
	The comparison concerns the method of pattern recognition, not the state of confirmation:
	
	\begin{table}[H]
		\centering
		\adjustbox{max width=\textwidth}{%
\begin{tabular}{p{3cm}p{4cm}p{4cm}p{3cm}}
			\toprule
			\textbf{Discovery} & \textbf{Pattern Recognition} & \textbf{Predictions} & \textbf{Confirmation} \\
			\midrule
			Periodic Table (1869) & Atomic weights and properties & Gallium, Germanium, Scandium & Experimentally confirmed \\
			Spectral Lines (1885) & Hydrogen lines & Rydberg formula for all series & Quantum mechanics \\
			Quark Model (1964) & Hadron masses & Eightfold Way & QCD theory \\
			\textbf{T0-Model (2025)} & \textbf{Lepton masses} & \textbf{Quarks} & \textbf{Experimental tests} \\
			\bottomrule
		\end{tabular}%
}
		\caption{Historical precedence of pattern physics}
		\label{tab:historical_precedence}
	\end{table}
	
	\section{From Energy Fields to Particle Masses}
	\label{sec:energy_fields_to_masses}
	
	\subsection{The Fundamental Challenge}
	\label{subsec:fundamental_challenge}
	
	A central result of the T0-Model is that particle masses can be calculated from geometric principles. While the Standard Model requires over 20 free parameters to describe particle masses, the T0-Model works with the geometric constant $\xigeom = \frac{4}{3} \times 10^{-4}$ and rational quantum numbers; the deviations for charged leptons and quarks are below 1.2\%.
	
	\begin{tcolorbox}[colback=green!5!white,colframe=green!75!black,title=Parameter Reduction for Masses]
		\textbf{Parameter Reduction:}
		\begin{itemize}
			\item \textbf{Standard Model}: 20+ free mass parameters (measured)
			\item \textbf{T0-Model}: 1 parameter $\xi$ plus rational quantum numbers (geometric)
			\item \textbf{Experimental Accuracy}: 99\% average agreement (charged leptons and quarks; neutrino part see addendum)
			\item \textbf{Theoretical Foundation}: Three-dimensional spatial geometry + QFT derivation
		\end{itemize}
	\end{tcolorbox}
	
	\subsection{Energy-Based Mass Concept}
	\label{subsec:energy_based_mass}
	
	In the T0 framework, what we traditionally call "mass" is interpreted as an expression of characteristic energy scales of field excitations:
	
	\begin{equation}
		\boxed{m_i \rightarrow E_{\text{char},i} \quad \text{(characteristic energy of particle type } i\text{)}}
		\label{eq:mass_to_energy}
	\end{equation}
	
	This reinterpretation treats mass and energy as different aspects of the same quantity.
	
	\section{Two Complementary Calculation Methods}
	\label{sec:two_calculation_methods}
	
	The T0-Model offers two mathematically equivalent but conceptually different approaches to calculating particle masses:
	
	\subsection{Method 1: Direct Geometric Resonance}
	\label{subsec:direct_geometric_method}
	
	\textbf{Conceptual Foundation:} Particles as resonances in the universal energy field
	
	The direct method treats particles as characteristic resonance modes of the energy field $\Efield$, analogous to standing wave patterns:
	
	\begin{equation}
		\text{Particle} = \text{Discrete resonance modes of } \Efield(x,t)
	\end{equation}
	
	\textbf{Three-Step Calculation Process:}
	
	\textbf{Step 1: Geometric Quantization}
	\begin{equation}
		\xi_i = \xi_0 \cdot f(n_i, l_i, j_i)
		\label{eq:geometric_quantization}
	\end{equation}

Here $f(n,l,j)$ is the geometric structure factor of the quantum numbers, not a frequency and not the base winding number $f = 1/\xi$.

	where:
	\begin{align}
		\xi_0 &= \frac{4}{3} \times 10^{-4} \quad \text{(geometric base parameter)} \\
		n_i, l_i, j_i &= \text{quantum numbers from 3D wave equation} \\
		f(n_i, l_i, j_i) &= \text{geometric function from spatial harmonics}
	\end{align}
	
	\textbf{Step 2: Resonance Frequencies}
	\begin{equation}
		\omega_i = \frac{c^2}{\xi_i \cdot r_{\text{char}}}
		\label{eq:resonance_frequencies}
	\end{equation}
	
	In natural units ($c = 1$):
	\begin{equation}
		\omega_i = \frac{1}{\xi_i}
	\end{equation}
	
	\textbf{Step 3: Mass Determination from Energy Conservation}
	\begin{equation}
		E_{\text{char},i} = \hbar \omega_i = \frac{\hbar}{\xi_i}
		\label{eq:energy_from_frequency}
	\end{equation}
	
	In natural units ($\hbar = 1$):
	\begin{equation}
		\boxed{E_{\text{char},i} = \frac{1}{\xi_i}}
		\label{eq:characteristic_energy_direct}
	\end{equation}
	
	\subsection{Method 2: Extended Yukawa Method}
	\label{subsec:extended_yukawa_method}
	
	\textbf{Conceptual Foundation:} Bridge to Standard Model formulation
	
	The extended Yukawa method maintains compatibility with Standard Model calculations while making Yukawa couplings geometrically determined rather than empirically fitted:
	
	\begin{equation}
		E_{\text{char},i} = y_i \cdot v
		\label{eq:yukawa_mass_formula}
	\end{equation}
	
	where $v = 246$ GeV is the Higgs vacuum expectation value (SM definition: $v\equiv(\sqrt{2}\,G_F)^{-1/2}$; no independent observable, Doc.~351/R112).
	
	\textbf{Geometric Yukawa Couplings:}
	\begin{equation}
		\boxed{y_i = r_i \cdot \left(\frac{4}{3} \times 10^{-4}\right)^{\pi_i}}
		\label{eq:geometric_yukawa}
	\end{equation}
	
	\textbf{Generation Hierarchy:}
	\begin{align}
		\text{1st Generation:} \quad &\pi_i = \frac{3}{2} \quad \text{(electron, up quark)} \\
		\text{2nd Generation:} \quad &\pi_i = 1 \quad \text{(muon, charm quark)} \\
		\text{3rd Generation:} \quad &\pi_i = \frac{2}{3} \quad \text{(tau, top quark)}
	\end{align}
	
	The coefficients $r_i$ are simple rational numbers determined by the geometric structure of each particle type.
	
	\section{Quantum Field Theoretical Derivation of the $\xi$-Constant}
	\label{sec:qft_derivation}
	
	\subsection{EFT Matching and Yukawa Coupling after EWSB}
	\label{subsec:eft_matching}
	
	After electroweak symmetry breaking, we have the Yukawa interaction:
	
	\begin{equation}
		\mathcal{L}_{\text{Yukawa}} \supset -\lambda_h \bar{\psi}\psi H, \quad \text{with} \quad H = \frac{v + h}{\sqrt{2}}
	\end{equation}
	
	After EWSB:
	\begin{equation}
		\mathcal{L} \supset -m \bar{\psi}\psi - y h \bar{\psi}\psi
	\end{equation}
	
	with the relations:
	\begin{equation}
		m = \frac{\lambda_h v}{\sqrt{2}} \quad \text{and} \quad y = \frac{\lambda_h}{\sqrt{2}}
	\end{equation}
	
	The local mass dependence on the physical Higgs field $h(x)$ leads to:
	
	\begin{equation}
		m(h) = m\left(1 + \frac{h}{v}\right) \quad \Rightarrow \quad \partial_\mu m = \frac{m}{v}\partial_\mu h
	\end{equation}
	
	\subsection{T0 Operators in Effective Field Theory}
	\label{subsec:t0_operators}
	
	In T0 theory, operators of the form appear:
	
	\begin{equation}
		O_T = \bar{\psi}\gamma^\mu\Gamma_\mu^{(T)}\psi
	\end{equation}
	
	with the characteristic time field coupling term:
	\begin{equation}
		\Gamma_\mu^{(T)} = \frac{\partial_\mu m}{m^2}
	\end{equation}
	
	Substituting the Higgs dependence:
	\begin{equation}
		\Gamma_\mu^{(T)} = \frac{\partial_\mu m}{m^2} = \frac{1}{mv}\partial_\mu h
	\end{equation}
	
	This suggests that a $\partial_\mu h$-coupled vector current is the UV origin.
	
	\subsection{1-Loop Matching Calculation}
	\label{subsec:one_loop_matching}
	
	The 1-loop amplitude for the T0 vertex yields:
	\begin{equation}
		F_V(0) = \frac{y^2}{16\pi^2}\left[\frac{1}{2} - \frac{1}{2}\ln\left(\frac{m_h^2}{\mu^2}\right) + \frac{r(r-\ln r-1)}{(r-1)^2}\right]
	\end{equation}
	
	For hierarchical masses ($m \ll m_h$), the constant term dominates:
	\begin{equation}
		F_V(0) \approx \frac{y^2}{32\pi^2}
	\end{equation}
	
	\subsection{$\xi$ Formula from Higgs Physics}
	\label{subsec:final_xi_formula}
	
	The EFT matching yields the relationship:
	\begin{equation}
		\boxed{\xi = \frac{\lambda_h^2 v^2}{16\pi^3 m_h^2}}
	\end{equation}
	
	With Standard Higgs parameters ($m_h = 125.1$ GeV, $v = 246.22$ GeV, $\lambda_h = m_h^2/(2v^2) \approx 0.129$):
	\begin{equation}
		\xi \approx 1.30 \times 10^{-4}
	\end{equation}
	
	This lies close to the geometric setting $\xi_0 = \frac{4}{3} \times 10^{-4} \approx 1.333 \times 10^{-4}$ (deviation $\approx -2.4\%$). The relation is a consistency check of $\xi$, not an exact derivation. For the origin of the Higgs route see Doc.~385.
	
	\section{Universal Particle Mass Systematics}
	\label{sec:universal_masses}
	
	\subsection{Revised Universal Table of Fermions}
	\label{subsec:universal_table}
	
\resizebox{\textwidth}{!}{%
	\begin{tabular}{|l|c|c|c|c|c|l|}
		\hline
		Fermion & Generation & Family & Spin & $r_f$ & Exponent $p_f$ & SM Correspondence \\
		\hline
		Electron Neutrino & 1 & 0 & 1/2 & $4/3$ & $5/2$ & Double $\xi$ \\
		Electron          & 1 & 0 & 1/2 & $4/3$  & $3/2$ & Lepton number \\
		Muon Neutrino     & 2 & 1 & 1/2 & $16/5$ & $3$   & Double $\xi$ \\
		Muon              & 2 & 1 & 1/2 & $16/5$ & $1$   & Lepton number \\
		Tau Neutrino      & 3 & 2 & 1/2 & $25/9$ & $25/9$ & Double $\xi$ \\
		Tau               & 3 & 2 & 1/2 & $25/9$  & $2/3$ & Lepton number \\
		\hline
		Up     & 1 & 0 & 1/2 & $6$          & $3/2$ & Color \\
		Down   & 1 & 0 & 1/2 & $\tfrac{25}{2}$ & $3/2$ & Color + Isospin \\
		Charm  & 2 & 1 & 1/2 & $2$$^*$          & $2/3$ & Color \\
		Strange& 2 & 1 & 1/2 & $\tfrac{26}{9}$ & $1$   & Color \\
		Top    & 3 & 2 & 1/2 & $\tfrac{1}{28}$ & $-1/3$ & Color \\
		Bottom & 3 & 2 & 1/2 & $\tfrac{3}{2}$  & $1/2$ & Color \\
		\hline
	\end{tabular}%
}
	
	\footnotetext{* Corrected from originally $8/9$ based on detailed numerical analysis}
	
	\section{Numerical Reconstruction}
	\label{sec:complete_numerical_reconstruction}
	
	The following analysis shows the explicit calculation of all fermions with both methods:
	
	\subsection{Foundations and Experimental Input Data}
	\label{subsec:foundations}
	
	\textbf{Fundamental Constants:}
	\begin{align}
		\xi_0 = \xi &= \frac{4}{3} \times 10^{-4} = 1.333333333\ldots \times 10^{-4} \\
		v &= 246 \text{ GeV}
	\end{align}
	
	\textbf{Experimental Masses (PDG-near values):}
	\begin{align}
		m_e^{\text{exp}} &= 0.0005109989461 \text{ GeV} \\
		m_\mu^{\text{exp}} &= 0.1056583745 \text{ GeV} \\
		m_\tau^{\text{exp}} &= 1.77686 \text{ GeV}
	\end{align}
	
	\subsection{Charged Leptons: Detailed Calculations}
	\label{subsec:charged_leptons_detailed}
	
	\textbf{Electron Mass Calculation:}
	
	\textit{Direct Method:}
	\begin{align}
		\xi_e &= \frac{4}{3} \times 10^{-4} \times f_e(1,0,1/2), \quad f_e = 1.485 \times 10^{7}\ \text{GeV}^{-1} \\
		&= 1980\ \text{GeV}^{-1} \\
		E_{e} &= \frac{1}{\xi_e} = 0.505 \text{ MeV}
	\end{align}
	
	\textit{Extended Yukawa Method:}
	\begin{align}
		r_e &= \frac{4}{3} \\
		y_e &= \frac{4}{3} \times \left(\frac{4}{3} \times 10^{-4}\right)^{3/2} = 2.053 \times 10^{-6} \\
		E_e &= y_e \times 246 \text{ GeV} = 0.505 \text{ MeV}
	\end{align}
	
	\textbf{Muon Mass Calculation:}
	
	\textit{Direct Method:}
	\begin{align}
		\xi_\mu &= \frac{4}{3} \times 10^{-4} \times f_\mu(2,1,1/2), \quad f_\mu = 7.146 \times 10^{4}\ \text{GeV}^{-1} \\
		&= 9.528\ \text{GeV}^{-1} \\
		E_{\mu} &= \frac{1}{\xi_\mu} = 104.96 \text{ MeV}
	\end{align}
	
	\textit{Extended Yukawa Method:}
	\begin{align}
		y_\mu &= \frac{16}{5} \times \left(\frac{4}{3} \times 10^{-4}\right)^1 = 4.267 \times 10^{-4} \\
		E_\mu &= y_\mu \times 246 \text{ GeV} = 104.96 \text{ MeV}
	\end{align}
	\textbf{Experiment:} $105.66 \text{ MeV}$ → Deviation $\approx 0.65\%$

	\textbf{Experiment electron:} $0.511 \text{ MeV}$ → Deviation $-1.18\%$

	The factors $f_i$ of the direct method follow from the fixed quantum numbers via Eq.~\ref{eq:transformation_bridge}: $f_e = 1/(r_e\,\xi^{5/2}\,v) = 1.485 \times 10^7$~GeV$^{-1}$ and $f_\mu = 1/(r_\mu\,\xi^{2}\,v) = 7.146 \times 10^4$~GeV$^{-1}$. They must not be equated with the Yukawa prefactors $r_i$. With these the direct method gives $1/(\xi f_e) = 0.505$~MeV and $1/(\xi f_\mu) = 104.96$~MeV, the same values as the Yukawa method.
	
	\subsection{Neutrino Treatment}
	\label{sec:complete_neutrino_treatment}
	
	\begin{neutrino}[Neutrino Approach]
		This section proposes a geometric treatment of neutrino masses via a \textbf{double $\xi$-suppression}. According to R109, this neutrino part is not tenable (see the addendum at the end); for neutrinos, Doc.~340 is authoritative.
	\end{neutrino}
	
	\subsection{Neutrino Quantum Numbers}
	\label{subsec:neutrino_quantum_numbers}
	
	Neutrinos follow the same quantum number structure as other fermions, but with a modification due to their weak interaction nature:
	
	\begin{table}[H]
		\centering
		\adjustbox{max width=\textwidth}{%
\begin{tabular}{lcccc}
			\toprule
			\textbf{Neutrino} & \textbf{n} & \textbf{l} & \textbf{j} & \textbf{Suppression} \\
			\midrule
			$\nu_e$ & 1 & 0 & 1/2 & Double $\xi$ \\
			$\nu_\mu$ & 2 & 1 & 1/2 & Double $\xi$ \\
			$\nu_\tau$ & 3 & 2 & 1/2 & Double $\xi$ \\
			\bottomrule
		\end{tabular}%
}
		\caption{Neutrino quantum numbers with characteristic double $\xi$-suppression}
		\label{tab:neutrino_quantum_numbers}
	\end{table}
	
	\subsection{Double $\xi$-Suppression Mechanism}
	\label{subsec:double_xi_suppression}
	
	The approach assumes that neutrinos experience an additional geometric suppression factor:
	
	\begin{equation}
		f(n_{\nu_i}, l_{\nu_i}, j_{\nu_i}) = f(n_i, l_i, j_i)_{\text{Lepton}} \times \xi
		\label{eq:neutrino_suppression}
	\end{equation}
	
	\textbf{Neutrino Mass Calculations:}
	
	\textbf{Electron Neutrino:}
	\begin{align}
		\xi_{\nu_e} &= \frac{4}{3} \times 10^{-4} \times 1 \times \frac{4}{3} \times 10^{-4} = \frac{16}{9} \times 10^{-8} \\
		E_{\nu_e} &= \frac{1}{\xi_{\nu_e}} = 9.1 \text{ meV}
	\end{align}
	
	\textbf{Muon Neutrino:}
	\begin{align}
		\xi_{\nu_\mu} &= \frac{4}{3} \times 10^{-4} \times \frac{16}{5} \times \frac{4}{3} \times 10^{-4} = \frac{256}{45} \times 10^{-8} \\
		E_{\nu_\mu} &= \frac{1}{\xi_{\nu_\mu}} = 1.9 \text{ meV}
	\end{align}
	
	\textbf{Tau Neutrino:}
	\begin{align}
		\xi_{\nu_\tau} &= \frac{4}{3} \times 10^{-4} \times \frac{25}{9} \times \frac{4}{3} \times 10^{-4} = \frac{400}{81} \times 10^{-8} \\
		E_{\nu_\tau} &= \frac{1}{\xi_{\nu_\tau}} = 18.0 \text{ meV}
	\end{align}
	
	\section{Quark Analysis with Both Methods}
	\label{sec:quark_analysis}
	
	\subsection{Explicit Calculations of Quark Masses}
	\label{subsec:quark_calculations}
	
	We use $\xi=\tfrac{4}{3}\times10^{-4}$ and $v=246\ \mathrm{GeV}$ (SM definition from $G_F$, Doc.~351/R112).
	For the Yukawa representation:
	\[
	y_i = r_i\,\xi^{p_i},\qquad m_i^{\text{pred}}=y_i\,v.
	\]
	For the direct geometric representation:
	\[
	f_i=\frac{1}{\xi\, m_i^{\text{exp}}},\qquad m_i^{\text{exp}}=\frac{1}{\xi\, f_i}.
	\]
	
	\begin{table}[H]
		\centering
			\resizebox{\textwidth}{!}{%
		\begin{tabular}{lcccccc}
			\toprule
			Quark & $p_i$ & $r_i$ (corr.) & $m_i^{\text{pred}}$ & $m_i^{\text{exp}}$ & rel.\ error & Remark\\
			& & & (GeV) & (GeV) & (\%) & \\
			\midrule
			Up     & $3/2$ & $6$        & $2.272\times10^{-3}$ & $2.27\times10^{-3}$ & $+0.11$ & OK \\
			Down   & $3/2$ & $25/2$     & $4.734\times10^{-3}$ & $4.72\times10^{-3}$ & $+0.30$ & OK \\
			Strange& $1$   & $26/9$        & $9.476\times10^{-2}$  & $9.50\times10^{-2}$  & $-0.26$ & OK\\
			Charm  & $2/3$ & $2$      & $1.284\times10^{0}$  & $1.28$              & $+0.32$ & Corrected\\
			Bottom & $1/2$ & $3/2$      & $4.261\times10^{0}$   & $4.26$              & $+0.02$ & OK \\
			Top    & $-1/3$& $1/28$     & $1.7198\times10^{2}$  & $171$               & $+0.57$ & OK \\
			\bottomrule
		\end{tabular}%
	}
		\caption{Yukawa predictions with corrected $r_i,p_i$ and comparison with reference masses.}
		\label{tab:quark_yukawa_predictions}
	\end{table}
	
	\subsection{Correction for the Charm Quark}
	\label{subsec:charm_correction}
	
	The originally specified value $r_c=8/9$ does not reproduce the reference mass $m_c=1.28\ \mathrm{GeV}$. The required value is:
	\[
	r_c^{\text{required}}=\frac{m_c^{\text{exp}}}{v\,\xi^{2/3}}\approx 1.994 \approx 2.
	\]
	
	Therefore, $r_c \approx 2$ was used in the corrected universal table. With $r_c = 2$ exactly, $m_c = 1.284$~GeV ($+0.32\,\%$).

		
	\section{Experimental Comparison}
	\label{sec:comprehensive_validation}
	
	\subsection{Accuracy Analysis}
	\label{subsec:complete_accuracy}
	
	Comparison of the T0 values with the experimental masses:
	
	\begin{table}[H]
		\centering
			\resizebox{\textwidth}{!}{%
		\begin{tabular}{lcccc}
			\toprule
			\textbf{Particle} & \textbf{T0 Prediction} & \textbf{Experiment} & \textbf{Accuracy} & \textbf{Type} \\
			\midrule
			\multicolumn{5}{c}{\textit{Charged Leptons}} \\
			\midrule
			Electron & 0.505 MeV & 0.511 MeV & 98.82\% & Lepton \\
			Muon & 104.96 MeV & 105.66 MeV & 99.35\% & Lepton \\
			Tau & 1783.4 MeV & 1776.86 MeV & 99.63\% & Lepton \\
			\midrule
			\multicolumn{5}{c}{\textit{Neutrinos}} \\
			\midrule
			$\nu_e$ & 9.1 meV & $< 450$ meV & Compatible & Neutrino \\
			$\nu_\mu$ & 1.9 meV & $< 180$ keV & Compatible & Neutrino \\
			$\nu_\tau$ & 18.0 meV & $< 18$ MeV & Compatible & Neutrino \\
			\midrule
			\multicolumn{5}{c}{\textit{Quarks}} \\
			\midrule
			Up quark & 2.272 MeV & 2.27 MeV & 99.89\% & Quark \\
			Down quark & 4.734 MeV & 4.72 MeV & 99.70\% & Quark \\
			Strange quark & 94.76 MeV & 95.0 MeV & 99.74\% & Quark \\
			Charm quark & 1.284 GeV & 1.28 GeV & 99.68\% & Quark \\
			Bottom quark & 4.261 GeV & 4.26 GeV & 99.98\% & Quark \\
			Top quark & 171.99 GeV & 171 GeV & 99.43\% & Quark \\
			\midrule
			\textbf{Average} & & & \textbf{99.6\%} & \textbf{All Fermions} \\
			\bottomrule
		\end{tabular}%
	}
		\caption{Comparison of T0 values with experimental masses}
		\label{tab:complete_validation}
	\end{table}
	
	
	\begin{keyresult}[Result for Charged Leptons and Quarks]
		The T0-Model achieves 99.6\% average accuracy across the charged leptons and quarks, with the single parameter $\xi$ and fixed rational quantum numbers. The neutrino rows are given only as compatibility with bounds; the neutrino part is not tenable according to R109 (see addendum).
	\end{keyresult}
	
	\section{Predictive Power of the Established System}
	\label{sec:predictive_power}
	
	\subsection{New Particle Generations}
	\label{subsec:new_generations}
	
	Extrapolating the established patterns to a 4th generation gives:
	
	\textbf{4th Generation (extrapolated):}
	\begin{align}
		n &= 4, \quad \pi_4 = \frac{1}{2}, \quad r_4 \approx 2.0 \\
		m_{\text{4thGen}} &= r_4 \times \xi^{1/2} \times v \approx 5.7 \text{ GeV}
	\end{align}

	The calculation gives $5.68$~GeV. A charged lepton of a 4th generation is excluded below about $100$~GeV by LEP, and the $Z$ width allows only three light neutrinos; this extrapolation is therefore not a viable prediction.
	
	\subsection{Quark Sector Extrapolation}
	\label{subsec:quark_extrapolation}
	
	The lepton patterns can be transferred to quarks:
	
	\begin{table}[H]
		\centering
		\adjustbox{max width=\textwidth}{%
\begin{tabular}{lcccc}
			\toprule
			\textbf{Quark} & \textbf{Generation} & \textbf{$r_i$} & \textbf{$\pi_i$} & \textbf{Prediction} \\
			\midrule
			Up & 1 & 6 & 3/2 & 2.3 MeV \\
			Down & 1 & 12.5 & 3/2 & 4.7 MeV \\
			Charm & 2 & 2.0 & 2/3 & 1.3 GeV \\
			Strange & 2 & 2.89 & 1 & 95 MeV \\
			Top & 3 & 0.036 & -1/3 & 173 GeV \\
			Bottom & 3 & 1.5 & 1/2 & 4.3 GeV \\
			\bottomrule
		\end{tabular}%
}
		\caption{Quark predictions from established patterns}
		\label{tab:quark_predictions}
	\end{table}
	
	\section{Corrected Interpretation of Mathematical Equivalence}
	\label{sec:corrected_interpretation}
	
	\begin{keypoint}[Meaning of the Equivalence]
		The mathematical equivalence of both methods is \textbf{given by definition} when the parameters ($r_i$ or $f_i$) are determined from the same experimental masses. The equivalence is not proof of the theory, but a consistency property of the mathematical structure.
	\end{keypoint}
	
	\subsection{Transformation Relationship as Bridge}
	\label{subsec:transformation_relationship}
	
	The relationship:
	\begin{equation}
		f_i = \frac{1}{r_i \, \xi^{\pi_i} \, v \, \xi_0}
		\label{eq:transformation_bridge}
	\end{equation}
	
	connects both methods mathematically. If $r_i$ is determined from experimental masses, $f_i$ follows automatically and vice versa.
	
	\begin{table}[H]
		\centering
		\adjustbox{max width=\textwidth}{%
\begin{tabular}{lcccc}
			\toprule
			\textbf{Particle} & \textbf{$m^{\text{exp}}$ (GeV)} & \textbf{$r_i$ (Yukawa)} & \textbf{$f_i$ (direct)} & \textbf{Accuracy} \\
			\midrule
			Electron & 0.000511 & 4/3 & $1.485 \times 10^{7}$ & $98.82\%$ \\
			Muon & 0.10566 & 16/5 & $7.146 \times 10^{4}$ & $99.34\%$ \\
			Tau & 1.77686 & 25/9 & $4.205 \times 10^{3}$ & $99.63\%$ \\
			\midrule
			$\nu_e$ & 9.1 $\times 10^{-6}$ & 1.349 & $8.235 \times 10^{10}$ & Prediction \\
			$\nu_\mu$ & 1.9 $\times 10^{-6}$ & 3.221 & $3.947 \times 10^{11}$ & Prediction \\
			$\nu_\tau$ & 18.0 $\times 10^{-6}$ & 2.778 & $3.989 \times 10^{10}$ & Prediction \\
			\bottomrule
		\end{tabular}%
}
		\caption{Numerical equivalence of both T0 methods for all leptons}
		\label{tab:numerical_equivalence_complete}
	\end{table}

	For the charged leptons the fixed quantum numbers $4/3$, $16/5$, $25/9$ are given with the $f_i$ following from them via Eq.~\ref{eq:transformation_bridge} (in GeV$^{-1}$); both methods give the same masses and hence the same accuracy. The neutrino rows belong to the neutrino part, which is not tenable (see addendum, R109).
	
	\section{Experimental Predictions and Precision Tests}
	\label{sec:experimental_predictions}
	
	\subsection{Modified QED Vertex Corrections}
	\label{subsec:qed_corrections}
	
	T0 theory predicts modified Feynman rules:
	\begin{align}
		\text{Time field vertex:} \quad &-i\gamma^\mu\Gamma_\mu^{(T)} = i\gamma^\mu\frac{\partial_\mu m}{m^2} \\
		\text{Modified fermion propagator:} \quad &S_F^{(T0)}(p) = S_F(p) \cdot \left[1 + \frac{\beta}{p^2}\right]
	\end{align}
	
	\subsection{Neutrino Validation}
	\label{subsec:neutrino_validation}
	
	The T0 neutrino values are compatible with the experimental bounds listed here (on the tenability of the neutrino part, see addendum, R109):
	
	\begin{table}[H]
		\centering
		\adjustbox{max width=\textwidth}{%
\begin{tabular}{lccc}
			\toprule
			\textbf{Parameter} & \textbf{T0 Prediction} & \textbf{Experimental Limit} & \textbf{Status} \\
			\midrule
			$m_{\nu_e}$ & 9.1 meV & $< 450$ meV (KATRIN) & Satisfied \\
			$m_{\nu_\mu}$ & 1.9 meV & $< 180$ keV (indirect) & Satisfied \\
			$m_{\nu_\tau}$ & 18.0 meV & $< 18$ MeV (indirect) & Satisfied \\
			$\sum m_\nu$ & 29.8 meV & $< 60$ meV (Cosmology 2024) & Satisfied \\
			\bottomrule
		\end{tabular}%
}
		\caption{T0 neutrino predictions vs. experimental constraints}
		\label{tab:neutrino_validation}
	\end{table}
	
	\begin{important}[Neutrino Mass Hierarchy]
		The T0-Model predicts \textbf{normal ordering}: $m_{\nu_\mu} < m_{\nu_e} < m_{\nu_\tau}$, which is consistent with current oscillation data preferences.
	\end{important}
	
	\section{Methodological Foundation}
	\label{sec:scientific_legitimacy}
	
	\subsection{Reversibility of the Established System}
	\label{subsec:reversibility}
	
	After the establishment phase, the T0 system can make predictions:
	
	\textbf{Established Lepton Patterns:}
	\begin{align}
		\text{1st Generation (n=1):} \quad &\pi_i = \frac{3}{2}, \quad r_e = \frac{4}{3} \approx 1.333 \\
		\text{2nd Generation (n=2):} \quad &\pi_i = 1, \quad r_\mu = \frac{16}{5} = 3.2 \\
		\text{3rd Generation (n=3):} \quad &\pi_i = \frac{2}{3}, \quad r_\tau = \frac{25}{9} \approx 2.778
	\end{align}
		
	\subsection{Experimental Testability}
	\label{subsec:experimental_testability}
	
	The T0 predictions are experimentally falsifiable:
	
	\begin{enumerate}
		\item \textbf{Precision Measurements:} Refinement of $r_i$ parameters
		\item \textbf{Neutrino Tests:} Direct neutrino mass measurements
	\end{enumerate}
	
	Methodologically, the following speaks for the T0 procedure:
	
	\begin{enumerate}
		\item \textbf{Systematic Structure:} All parameters follow recognizable patterns
		\item \textbf{Predictive Power:} After establishment, new particles become predictable
		\item \textbf{Experimental Testability:} Predictions are falsifiable
		\item \textbf{QFT Foundation:} Quantum field theoretical derivation of the $\xi$-constant
		\item \textbf{Historical Models:} The methodology of pattern physics has been successful several times
	\end{enumerate}
	
	\section{Parameter Economy and Universal Structure}
	\label{sec:parameter_free_nature}
	
	\begin{important}[No Adjustable Parameters]
		The scale of the T0 coefficients is determined by $\xi$, which can also be checked against Higgs parameters (deviation from $4/3 \times 10^{-4}$ about $-2.4$\%):
		\begin{equation}
			\xi = \frac{\lambda_h^2 v^2}{16\pi^3 m_h^2} \approx 1.30 \times 10^{-4}
		\end{equation}
		The relation is a consistency check, not a derivation (see Section~\ref{subsec:final_xi_formula}).
		Together with the rational quantum numbers $(r_i, p_i)$, no continuously adjustable parameter remains; $v$ serves as a conversion anchor.
	\end{important}
	
	\subsection{Universal Quantum Number Table}
	\label{subsec:universal_quantum_table}
	
	\begin{table}[H]
		\centering
		\adjustbox{max width=\textwidth}{%
\begin{tabular}{lcccccc}
			\toprule
			\textbf{Particle} & \textbf{n} & \textbf{l} & \textbf{j} & \textbf{$r_i$} & \textbf{$p_i$} & \textbf{Special} \\
			\midrule
			\multicolumn{7}{c}{\textit{Charged Leptons}} \\
			\midrule
			Electron & 1 & 0 & 1/2 & 4/3 & 3/2 & -- \\
			Muon & 2 & 1 & 1/2 & 16/5 & 1 & -- \\
			Tau & 3 & 2 & 1/2 & 25/9 & 2/3 & -- \\
			\midrule
			\multicolumn{7}{c}{\textit{Neutrinos}} \\
			\midrule
			$\nu_e$ & 1 & 0 & 1/2 & 4/3 & 5/2 & Double $\xi$ \\
			$\nu_\mu$ & 2 & 1 & 1/2 & 16/5 & 3 & Double $\xi$ \\
			$\nu_\tau$ & 3 & 2 & 1/2 & 25/9 & 25/9 & Double $\xi$ \\
			\midrule
			\multicolumn{7}{c}{\textit{Quarks}} \\
			\midrule
			Up & 1 & 0 & 1/2 & 6 & 3/2 & Color \\
			Down & 1 & 0 & 1/2 & 25/2 & 3/2 & Color + Isospin \\
			Charm & 2 & 1 & 1/2 & 2 & 2/3 & Color \\
			Strange & 2 & 1 & 1/2 & 26/9 & 1 & Color \\
			Top & 3 & 2 & 1/2 & 1/28 & -1/3 & Color \\
			Bottom & 3 & 2 & 1/2 & 3/2 & 1/2 & Color \\
			\bottomrule
		\end{tabular}%
}
		\caption{Universal quantum number table for all fermions}
		\label{tab:universal_quantum_numbers}
	\end{table}
	
	
\subsubsection{Winding numbers and topological assignment}

The rational quantum numbers $(r_i, p_i)$ are not arbitrary.
They follow from the winding numbers $(n_\theta, n_\phi)$ of
the 4D torus $T^4$ via $r_i = n_\phi^2 / n_\theta$ (Doc.~317,
August 2026):

\begin{center}
\begin{tabular}{lcccccc}
\toprule
Lepton & Gen. & $n_\theta$ & $n_\phi$ & $r_i$ & $p_i$ &
  KSAU knot \\
\midrule
$e$       & 1 & 3 & 2 & $4/3$  & $3/2$ & $3_1$ (trefoil) \\
$\mu$     & 2 & 5 & 4 & $16/5$ & $1$   & $6_3$ (amphicheiral) \\
$\tau$    & 3 & 9 & 5 & $25/9$ & $2/3$ & $7_1$ (torus knot) \\
$\nu_e$   & 1 & 3 & 2 & $4/3$  & $5/2$ & --- \\
$\nu_\mu$ & 2 & 5 & 4 & $16/5$ & $3$   & --- \\
$\nu_\tau$& 3 & 9 & 5 & $25/9$ & $25/9$& --- \\
\bottomrule
\end{tabular}
\end{center}

The exponent step $\Delta p = 1/3$ corresponds to the modulation
frequency of the torus winding. The knot assignment (KSAU,
Zenodo 2026) is \textbf{[S]}; the FFGFT quantum numbers are
\textbf{[K]}. Details in Doc.~317 (KSAU-FFGFT synthesis).

\section{Critical Evaluation and Limitations}
	\label{sec:critical_evaluation}
	
	\subsection{Theoretical Open Questions}
	\label{subsec:open_questions}
	
	\begin{enumerate}
		\item \textbf{Number of Generations:} Why exactly three generations?
		\item \textbf{Hierarchy Problem:} Connection between different energy scales
		\item \textbf{CP Violation:} Incorporation of CKM and PMNS mixing matrices
	\end{enumerate}
	
	\section{Concluding Assessment}
	\label{sec:concluding_assessment}
	
	\subsection{Scientific Status}
	\label{subsec:scientific_status}
	
	The T0-Model offers a systematic description of particle masses. The combination of:
	
	\begin{itemize}
		\item \textbf{High numerical accuracy} (99.6\% across charged leptons and quarks)
		\item \textbf{Parameter economy} (one parameter $\xi$, rational quantum numbers)
		\item \textbf{Broad coverage} (charged leptons and quarks; neutrinos see addendum)
		\item \textbf{QFT consistency} (1-loop derivation of the $\xi$-constant)
		\item \textbf{Experimental testability} (specific falsifiable predictions)
	\end{itemize}
	
	makes the model a testable proposal.
	
	\subsection{Significance for Fundamental Physics}
	\label{subsec:significance_physics}
	
	If experimentally confirmed, the T0-Model would considerably change our understanding of particle physics:
	
	\begin{enumerate}
		\item \textbf{Geometric Interpretation:} Particle masses as manifestations of 3D spatial geometry
		\item \textbf{Unification:} All fermions follow the same universal structure
		\item \textbf{Predictive Power:} New particles become predictable from established patterns
		\item \textbf{Simplicity:} Strong simplification of complex phenomena
	\end{enumerate}
	
	The T0-Model is an example showing that a unified description does not necessarily require greater complexity; the basic structure might be simple.
	
	\begin{thebibliography}{99}
		\bibitem{pascher_t0_energy_2025}
		Pascher, J. (2025). \textit{Das T0-Modell (Planck-referenziert): Eine Reformulierung der Physik}. 
		
		\bibitem{pascher_derivation_2025}
		Pascher, J. (2025). \textit{Feldtheoretische Ableitung des $\beta_T$-Parameters in natürlichen Einheiten ($\hbar = c = 1$)}. 
		
		\bibitem{pascher_qft_2025}
		Pascher, J. (2025). \textit{Vollständige Herleitung der Higgs-Masse und Wilson-Koeffizienten}. T0-Theory Project Documentation.
		
		\bibitem{pascher_units_2025}  
		Pascher, J. (2025). \textit{Natürliche Einheitensysteme: Universelle Energiekonversion und fundamentale Längenskala-Hierarchie}. 
		
		\bibitem{katrin_2024}
		KATRIN Collaboration. (2024). \textit{Direct neutrino mass measurement based on 259 days of KATRIN data}. arXiv:2406.13516.
		
		\bibitem{nufit_2024}
		Esteban, I., et al. (2024). \textit{NuFit-6.0: Updated global analysis of three-flavor neutrino oscillations}. J. High Energy Phys. 12, 216.
		
		\bibitem{cosmology_2024}
		Planck Collaboration. (2020). \textit{Planck 2018 results. VI. Cosmological parameters}. Astron. Astrophys. 641, A6.
		
		\bibitem{gell_mann_1964}
		Gell-Mann, M. (1964). \textit{A schematic model of baryons and mesons}. Physics Letters, 8(3), 214--215.
		
		\bibitem{mendeleev_1869}
		Mendeleev, D. (1869). \textit{Über die Beziehungen der Eigenschaften zu den Atomgewichten der Elemente}. Zeitschrift für Chemie, 12, 405--406.
		
		\bibitem{muon_g2_2023}
		Muon g-2 Collaboration. (2023). \textit{Measurement of the positive muon anomalous magnetic moment to 0.20 ppm}. Phys. Rev. Lett. 131, 161802.
		
	\end{thebibliography}

% ============================================================
% Addendum 7 September 2026 — Status clarification (R105, R109)
% Source: Doc. 190-Archive-2, declared 1 September 2026
% Source document remains unchanged (append-only, cf. R50).
% ============================================================

\section*{Addendum: Status Clarification (7~September 2026)}

\paragraph{R105: Retroactive marker certification (1~September 2026).}
This document was created before the marker system was introduced (Doc.~A010)
and carries no inline markers. Its derivations are internally consistent and
serve as chain links for the Galois programme (Docs.~336, 338, 339).

Retroactive status for this document:
Winding numbers $r_i$, exponents $p_i$, and the Yukawa method
$m_i = r_i\,\xi^{p_i}\,v$ for charged leptons and quarks: \textbf{[K]}.

Dependent documents may treat these results as a secured chain link
via reference to R105, without modifying this document.

\paragraph{R109: Restriction --- neutrino part not viable \textbf{[X]} (1~September 2026).}
The ``direct method'' $E_i = 1/\xi_i$ for neutrinos is not a calculation:
$1/(\xi r_i)$ gives the ratio $1:0.42$ for $e:\mu$ (inverse hierarchy,
dimensionless); the MeV values alongside are measurements.
The neutrino masses $9.1/1.9/18.0$~meV follow from three mutually
incompatible prescriptions; the exponent $p_{\nu_\tau}=25/9$ is $r_\tau$,
erroneously copied into the exponent column.

\textbf{Valid:} the Yukawa method $m_i=r_i\xi^{p_i}v$ for charged leptons
and quarks ($<1.2\,\%$ deviation). The \textbf{[K]} certification from R105
applies to this document \emph{only} for $r_i$, $p_i$, and the Yukawa method
--- not for the neutrino part and not for the claimed equivalence of both methods.

\textbf{Clarification:} Both methods are exactly
equivalent via $f_i=1/(r_i\,\xi^{p_i+1}v)$ (Eq.~\ref{eq:transformation_bridge}) and
give the same masses with fixed $r_i$, $p_i$.

\textbf{Authoritative for neutrinos:} Doc.~340 (R103, R108).
